\documentclass[aspectratio=169]{beamer}

% =====================================================================
% ART DECK 08 -- THE HAND-COLOURED ENGRAVING.
% The pen-and-ink plates of deck 06, but watercolour washes bloom in
% plate by plate: wash opacity is driven by a per-plate saturation knob
% (\setsat). Frontispiece = pure ink; the Summit = fully tinted. The
% black linework never changes -- only the colour fills in, as on an
% antique hand-coloured print. Two pdflatex passes (overlay).
% =====================================================================

\usepackage[T1]{fontenc}
\usepackage{ebgaramond}
\renewcommand{\familydefault}{\rmdefault}
\usepackage{tikz}
\usetikzlibrary{decorations.pathmorphing, calc, arrows.meta, shapes.misc}

\usetheme{default}
\usefonttheme{professionalfonts}
\setbeamertemplate{navigation symbols}{}
\setbeamertemplate{footline}{}

\definecolor{paper}{HTML}{FAF7EF}
\definecolor{ink}{HTML}{211E18}
\definecolor{inkl}{HTML}{6B6557}
% muted watercolour pigments
\definecolor{wsky}{HTML}{8FB0C9}
\definecolor{wwater}{HTML}{6E97B5}
\definecolor{wgreen}{HTML}{7E9347}
\definecolor{wgreen2}{HTML}{5C7A3E}
\definecolor{wamber}{HTML}{DBA23E}
\definecolor{wochre}{HTML}{B5824A}
\definecolor{wslate}{HTML}{7A8694}
\definecolor{wred}{HTML}{C24B43}
\definecolor{wviolet}{HTML}{8A6BA8}
\definecolor{wteal}{HTML}{3E9C8F}
\definecolor{wgold}{HTML}{CCA63E}

\setbeamercolor{background canvas}{bg=paper}
\setbeamercolor{normal text}{fg=ink}

\pgfmathsetseed{206}

\tikzset{
  pen/.style={draw=ink, line width=0.5pt, line cap=round, line join=round},
  penf/.style={draw=ink, line width=0.35pt, line cap=round, line join=round},
  pen2/.style={draw=ink, line width=0.9pt, line cap=round, line join=round},
  wob/.style={decorate, decoration={random steps, segment length=7pt, amplitude=0.4pt}},
  % paper-backed label so text never collides with the artwork
  lab/.style={fill=paper, fill opacity=0.9, text opacity=1, rounded corners=2pt, inner sep=3pt},
}

% saturation knob -> wash opacity \wo (0..1)
\newcommand{\setsat}[1]{\pgfmathsetmacro{\wo}{#1/100}}
\setsat{0}
% wash: fill the current clip region with pigment at wash strength (USE in a clip scope)
\newcommand{\wfill}[2]{\pgfmathsetmacro{\woo}{\wo*#2}\fill[#1, opacity=\woo] (-3,-3) rectangle (19,12);}

\newcommand{\art}[1]{%
  \begin{frame}[plain]
  \begin{tikzpicture}[remember picture, overlay,
      shift={(current page.south west)}, x=1cm, y=1cm]
    \fill[paper] (0,0) rectangle (16,9);
    \foreach \foxk in {1,...,40}{\pgfmathsetmacro{\fx}{rnd*16}\pgfmathsetmacro{\fy}{rnd*9}
       \fill[inkl, opacity=0.06] (\fx,\fy) circle (0.01);}
    #1
  \end{tikzpicture}
  \end{frame}}

\newcommand{\hatch}[4]{%
  \begin{scope}[rotate=#1]
    \pgfmathtruncatemacro{\hn}{52/#2}
    \foreach \hk in {0,...,\hn}{
      \pgfmathsetmacro{\hy}{-26+\hk*#2}
      \draw[#4, line width=#3] (-30,\hy) -- (30,\hy);}
  \end{scope}}

\newcommand{\stipple}[5]{%
  \foreach \si in {1,...,#5}{
    \pgfmathsetmacro{\sx}{#1+rnd*(#3-#1)}
    \pgfmathsetmacro{\sy}{#2+rnd*(#4-#2)}
    \fill[ink] (\sx,\sy) circle (0.012);}}

\newcommand{\platetitle}[1]{%
  \node[ink] at (8,8.45) {\scshape\fontsize{20}{22}\selectfont #1};
  \draw[penf] (4.2,8.05) -- (11.8,8.05);
  \draw[penf] (5.0,7.95) -- (11.0,7.95);}
\newcommand{\platecap}[1]{\node[ink] at (8,0.7) {\itshape #1};}

\newcommand{\pine}[2]{%
  \begin{scope}[shift={#1}, scale=#2]
    \draw[pen2] (0,-0.1) -- (0,0.35);
    \foreach \ty in {0.30,0.52,...,2.35}{
       \pgfmathsetmacro{\w}{1.05*(2.65-\ty)/2.65}
       \draw[pen] (0,\ty+0.22) -- (-\w,\ty);
       \draw[pen] (0,\ty+0.22) -- (\w,\ty);
       \foreach \k in {0.25,0.5,0.75}{
          \draw[penf] (-\w*\k,{\ty+0.22-0.22*\k}) -- (-\w*\k-0.14,{\ty+0.05-0.22*\k});
          \draw[penf] (\w*\k,{\ty+0.22-0.22*\k}) -- (\w*\k+0.14,{\ty+0.05-0.22*\k});}}
    \draw[pen2] (0,2.30) -- (0,2.66);
  \end{scope}}

% engraved sphere with an optional colour wash {(c)}{r}{spacing}{pigment}
\newcommand{\orb}[4]{%
  \begin{scope}\clip #1 circle (#2);\wfill{#4}{0.6}\end{scope}
  \begin{scope}
    \clip #1 circle (#2);
    \hatch{30}{#3}{0.4pt}{ink}\hatch{-35}{#3}{0.4pt}{ink}
    \begin{scope}\clip #1 ++(0.18,-0.18) circle (#2);
       \hatch{75}{#3}{0.4pt}{ink}\end{scope}
  \end{scope}
  \fill[paper] ($#1+(-0.32*#2,0.34*#2)$) circle (0.16*#2);
  \draw[pen2] #1 circle (#2);}

\newcommand{\rock}[3]{%
  \begin{scope}\clip[shift={#1}] (0,0) ellipse (#2 and #3);\wfill{wochre}{0.5}\end{scope}
  \begin{scope}
    \clip[shift={#1}] (0,0) ellipse (#2 and #3);
    \clip[shift={#1}] (-#2,-#3) rectangle (#2,{-0.12*#3});
    \hatch{25}{0.085}{0.4pt}{ink}\hatch{-30}{0.11}{0.4pt}{ink}
  \end{scope}
  \draw[pen2, wob, shift={#1}] (0,0) ellipse (#2 and #3);}

\begin{document}

% =====================================================================
% 1. FRONTISPIECE  -- pure ink (colour = 0)
% =====================================================================
\art{
  \setsat{0}
  \draw[pen] (0.7,0.7) rectangle (15.3,8.3);
  \draw[penf] (0.95,0.95) rectangle (15.05,8.05);
  \foreach \cx/\cy in {0.7/0.7,15.3/0.7,0.7/8.3,15.3/8.3}{
     \draw[pen] (\cx,\cy)+(-0.18,0) -- ++(0.18,0);
     \draw[pen] (\cx,\cy)+(0,-0.18) -- ++(0,0.18);}
  \begin{scope}
    \clip (1.2,4.4) -- (3.6,6.4) -- (5.2,5.2) -- (7.4,6.9) -- (9.0,5.6)
       -- (11.6,7.0) -- (13.2,5.3) -- (14.8,6.2) -- (14.8,4.4) -- cycle;
    \hatch{58}{0.16}{0.35pt}{inkl}
  \end{scope}
  \draw[pen, inkl] (1.2,4.4) -- (3.6,6.4) -- (5.2,5.2) -- (7.4,6.9) -- (9.0,5.6)
       -- (11.6,7.0) -- (13.2,5.3) -- (14.8,6.2);
  \begin{scope}
    \clip (9.8,4.4) -- (11.8,6.1) -- (13.6,5.0) -- (15.05,5.7) -- (15.05,4.4) -- cycle;
    \hatch{50}{0.1}{0.4pt}{ink}
  \end{scope}
  \draw[pen] (9.8,4.4) -- (11.8,6.1) -- (13.6,5.0) -- (15.05,5.7);
  \pine{(11.9,5.3)}{0.5}\pine{(12.7,5.0)}{0.42}\pine{(13.5,4.9)}{0.46}
  \begin{scope}
    \clip (0.95,0.95) -- (5.6,0.95) -- (4.4,3.6) -- (0.95,4.3) -- cycle;
    \hatch{80}{0.07}{0.4pt}{ink}
  \end{scope}
  \pine{(2.5,3.4)}{1.5}
  \rock{(5.4,2.0)}{0.55}{0.4}\rock{(6.4,1.75)}{0.6}{0.42}
  \rock{(7.5,1.7)}{0.55}{0.38}\rock{(4.5,2.2)}{0.45}{0.34}
  \foreach \ry in {1.2,1.45,...,3.4}{
     \draw[penf] (7.8,\ry) .. controls (10.2,{\ry+0.07}) and (12.2,{\ry-0.05}) .. (14.6,{\ry+0.12});}
  \fill[paper, opacity=0.93] (3.0,2.4) rectangle (13.0,4.5);
  \node[ink] at (8,3.8) {\scshape\fontsize{40}{44}\selectfont Rotor in Rust};
  \node[ink] at (8,3.0) {\itshape a field study in fast, trustworthy verification};
  \node[inkl] at (8,2.65) {\footnotesize Jasmin Begi\'c \quad$\cdot$\quad Daniel Wassie};
}

% =====================================================================
% 2. THE HAYSTACK  (colour 7)
% =====================================================================
\art{
  \setsat{7}
  \platetitle{The Haystack}
  \begin{scope}\clip (1.6,1.6) rectangle (14.4,6.6);\wfill{wamber}{0.30}\end{scope}
  \begin{scope}\clip (1.6,1.6) rectangle (14.4,6.6);\stipple{1.6}{1.6}{14.4}{6.6}{2300}\end{scope}
  \draw[pen] (1.6,1.6) rectangle (14.4,6.6);
  \fill[paper] (10.4,4.6) circle (0.55);
  \draw[pen2] (10.4,4.6) circle (0.32);
  \foreach \a in {0,30,...,330}{\draw[penf] (10.4,4.6)+(\a:0.62) -- +(\a:0.92);}
  \node[ink, lab] at (10.4,3.4) {\scriptsize the one that crashes};
  \platecap{eight bytes of input --- $2^{64}$ of them --- and you may try only a handful}
}

% =====================================================================
% 3. ONE QUESTION  (colour 15)
% =====================================================================
\art{
  \setsat{15}
  \platetitle{One Question, Asked of All}
  \begin{scope}\clip (1.6,1.4) rectangle (14.4,3.0);\wfill{wgreen}{0.4}\end{scope}
  \draw[pen] (1.6,3.0) -- (14.4,3.0);
  \begin{scope}\clip (1.6,1.4) rectangle (14.4,3.0);
     \hatch{0}{0.16}{0.35pt}{inkl}\end{scope}
  \begin{scope}\clip (8,5.6) circle (0.7);\wfill{wamber}{0.7}\end{scope}
  \draw[pen2] (8,5.6) circle (0.7);
  \begin{scope}\clip (8,5.6) circle (0.7);
     \hatch{20}{0.12}{0.35pt}{inkl}\end{scope}
  % short rays crown the upper half of the disc
  \foreach \a in {18,36,...,162}{
     \draw[penf] (8,5.6)+(\a:0.85) -- +(\a:{1.1+0.4*abs(sin(\a*3))});}
  % light falls below: beams start at the lower rim (never crossing the disc)
  \foreach \x in {2.4,3.2,...,13.6}{
     \draw[penf, inkl] ($(8,5.6)!0.82cm!(\x,3.0)$) -- (\x,3.0);}
  \platecap{can it reach a bad state, for \emph{any} input, within $k$ steps? --- one query covers them all}
}

% =====================================================================
% 4. THE APPARATUS  (colour 26)
% =====================================================================
\art{
  \setsat{26}
  \platetitle{The Apparatus}
  \foreach \x/\lbl/\wc in {2.6/binary/wslate, 6.4/model/wamber,
        10.2/solver/wslate, 13.6/picture/wamber}{
     \begin{scope}\clip (\x-0.95,2.7) rectangle (\x+0.95,5.3);\wfill{\wc}{0.5}\end{scope}
     \draw[pen2] (\x-0.95,2.7) -- (\x-0.95,5.3);
     \draw[pen2] (\x+0.95,2.7) -- (\x+0.95,5.3);
     \draw[pen2] (\x,5.3) ellipse (0.95 and 0.32);
     \draw[pen] (\x,2.7) ellipse (0.95 and 0.32);
     \begin{scope}\clip (\x-0.95,2.7) rectangle (\x+0.95,5.3);
        \clip (\x+0.2,2.4) rectangle (\x+0.95,5.3);
        \hatch{90}{0.1}{0.4pt}{ink}\end{scope}
     \node[ink, lab] at (\x,4.0) {\small\scshape \lbl};}
  \foreach \x in {3.55,7.35,11.15}{
     \draw[pen2, -{Stealth[length=2.4mm]}] (\x+0.02,4.0) -- (\x+0.9,4.0);}
  \node[inkl] at (4.5,5.9) {\footnotesize\itshape rotor};
  \node[inkl] at (11.9,5.9){\footnotesize\itshape visualiser};
  \draw[penf,-{Stealth[length=2mm]}] (4.5,5.7) to[bend right=20] (5.6,5.5);
  \draw[penf,-{Stealth[length=2mm]}] (12.4,5.7) to[bend left=20] (13.4,5.5);
  \platecap{the model checker reasons over the real machine code --- exactly what runs}
}

% =====================================================================
% 5. THE MECHANISM  (colour 37)
% =====================================================================
\art{
  \setsat{37}
  \platetitle{The Mechanism}
  \def\cx{8}\def\cy{4.7}
  \begin{scope}\clip (\cx,\cy) circle (2.5);\wfill{wgold}{0.55}\end{scope}
  \foreach \i in {0,...,23}{\pgfmathsetmacro{\a}{\i*15}
     \ifnum\i=7
        \draw[pen2] (\cx,\cy)+(\a:2.5) -- +(\a:3.05);
        \draw[pen2] (\cx,\cy)+({\a-3}:2.5) -- +({\a-3}:2.95);
        \draw[pen2] (\cx,\cy)+({\a+3}:2.5) -- +({\a+3}:2.95);
     \else
        \fill[ink] (\cx,\cy)+({\a-3.5}:2.5) -- +({\a-3.5}:2.85)
           -- +({\a+3.5}:2.85) -- +({\a+3.5}:2.5) -- cycle;
     \fi}
  \begin{scope}\clip (\cx,\cy) circle (2.5);
     \hatch{35}{0.14}{0.4pt}{ink}\hatch{-38}{0.18}{0.4pt}{inkl}\end{scope}
  \draw[pen2] (\cx,\cy) circle (2.5);
  \fill[paper] (\cx,\cy) circle (1.1);
  \draw[pen2] (\cx,\cy) circle (1.1);
  \node[ink, align=center] at (\cx,\cy) {\small state\\ rules\\ \textsc{24 limits}};
  \node[ink] at (11.8,6.5) {\footnotesize one tooth $=$};
  \node[ink] at (11.9,6.1) {\footnotesize \itshape divide-by-zero};
  \draw[penf,-{Stealth[length=2mm]}] (11.4,6.4) to[bend left=10] (8.78,7.1);
  \platecap{a position, the legal moves, and the squares the machine must never reach}
}

% =====================================================================
% 6. THE REWRITE  (colour 47)
% =====================================================================
\art{
  \setsat{47}
  \platetitle{Step the First --- Order from Tangle}
  \begin{scope}\clip (3.6,4.4) circle (1.5);
     \foreach \i in {1,...,90}{\pgfmathsetmacro{\a}{rnd*360}\pgfmathsetmacro{\b}{rnd*360}
        \pgfmathsetmacro{\r}{0.3+rnd*1.3}\pgfmathsetmacro{\s}{0.3+rnd*1.3}
        \draw[penf] (3.6,4.4)+(\a:\r) .. controls (3.6,4.4) .. +(\b:\s);}
  \end{scope}
  \draw[pen2] (3.6,4.4) circle (1.5);
  \node[ink, align=center] at (3.6,2.4) {\footnotesize \texttt{rotor.c}\\ \itshape 14{,}000 lines, one file};
  \draw[pen2,-{Stealth[length=2.6mm]}] (5.4,4.4) -- (7.2,4.4);
  \node[ink] at (6.3,4.9) {\footnotesize\itshape rebuild};
  \draw[pen2] (9.0,2.4) rectangle (13.6,6.6);
  \foreach \y/\n in {6.0/btor2, 5.0/riscv, 4.0/machine, 3.0/model}{
     \begin{scope}\clip (9.2,\y-0.45) rectangle (13.4,\y+0.45);\wfill{wochre}{0.5}\end{scope}
     \draw[pen] (9.2,\y-0.45) rectangle (13.4,\y+0.45);
     \draw[pen] (12.95,\y) ellipse (0.1 and 0.06);
     \node[ink] at (11.0,\y) {\small\scshape \n};}
  \platecap{same work --- now in four clean parts you can actually read}
}

% =====================================================================
% 7. SWIFTNESS  (colour 57)
% =====================================================================
\art{
  \setsat{57}
  \platetitle{Step the Second --- Swiftness}
  \foreach \x/\top/\lbl in {4.4/0.85/{139 s}, 11.6/0.06/{0.1 s}}{
     \draw[pen2] (\x-1.1,6.2) -- (\x+1.1,6.2);
     \draw[pen2] (\x-1.1,2.4) -- (\x+1.1,2.4);
     \draw[pen2] (\x-1.0,6.2) -- (\x,4.3) -- (\x-1.0,2.4);
     \draw[pen2] (\x+1.0,6.2) -- (\x,4.3) -- (\x+1.0,2.4);
     \begin{scope}\clip (\x-1.0,6.2) -- (\x,4.3) -- (\x+1.0,6.2) -- cycle;
        \clip (\x-1.1,{4.3+\top*1.9}) rectangle (\x+1.1,6.2);
        \wfill{wamber}{0.7}
        \hatch{20}{0.07}{0.4pt}{ink}\hatch{-24}{0.09}{0.4pt}{ink}\end{scope}
     \begin{scope}\clip (\x-1.0,2.4) -- (\x,4.3) -- (\x+1.0,2.4) -- cycle;
        \clip (\x-1.1,2.4) rectangle (\x+1.1,{2.4+(1-\top)*1.7});
        \wfill{wamber}{0.7}
        \hatch{20}{0.07}{0.4pt}{ink}\hatch{-24}{0.09}{0.4pt}{ink}\end{scope}
     \node[ink] at (\x,1.7) {\small \lbl};}
  \node[ink] at (8,5.0) {\scshape\fontsize{30}{32}\selectfont $\approx$1000$\times$};
  \draw[penf] (6.7,4.5) -- (9.3,4.5);
  \platecap{139 seconds to a tenth of one --- and 428 megabytes down to twenty}
}

% =====================================================================
% 8. THE COUNT  (colour 66)
% =====================================================================
\art{
  \setsat{66}
  \platetitle{Do Not Trust --- Count}
  \begin{scope}\clip (3.0,1.6) rectangle (5.4,6.8);\wfill{wslate}{0.5}\end{scope}
  \draw[pen2] (3.0,1.6) rectangle (5.4,6.8);
  \foreach \ly in {1.9,2.2,...,6.5}{
     \foreach \g in {0,1,2,3}{\draw[penf] (3.25+\g*0.22,\ly) -- (3.25+\g*0.22,\ly+0.18);}
     \draw[penf] (3.18,\ly+0.09) -- (3.95,\ly+0.02);}
  \node[ink, align=center] at (4.2,1.0) {\footnotesize 11.7 billion comparisons};
  \node[inkl] at (4.2,7.2) {\footnotesize\itshape walk the whole list};
  \begin{scope}\clip (10.0,3.8) rectangle (13.4,5.6);\wfill{wamber}{0.5}\end{scope}
  \draw[pen2] (10.0,3.8) rectangle (13.4,5.6);
  \draw[pen] (10.0,5.2) -- (13.4,5.2);
  \fill[ink] (10.5,4.5) circle (0.05);
  \draw[penf] (10.8,4.5) -- (12.9,4.5);
  \node[ink, align=center] at (11.7,3.2) {\footnotesize 159 thousand probes};
  \node[inkl] at (11.7,6.0) {\footnotesize\itshape one look-up};
  \node[ink] at (7.7,4.6) {\itshape vs.};
  \platecap{the same question, asked the same way --- a list, against an index}
}

% =====================================================================
% 9. THREE CASTINGS  (colour 76)
% =====================================================================
\art{
  \setsat{76}
  \platetitle{Three Castings of One Engine}
  \orb{(3.4,5.0)}{1.25}{0.10}{wamber}
  \orb{(8.0,5.0)}{1.25}{0.15}{wviolet}
  \orb{(12.6,5.0)}{1.25}{0.26}{wteal}
  \node[ink] at (3.4,3.2) {\scshape original C};
  \node[ink] at (8.0,3.2) {\scshape C \& the trick};
  \node[ink] at (12.6,3.2){\scshape rust};
  \node[inkl, align=center] at (3.4,2.5) {\footnotesize 139 s $\cdot$ 428 MB\\ 11.7 B compares};
  \node[inkl, align=center] at (8.0,2.5) {\footnotesize 1.5 s $\cdot$ 95$\times$\\ \itshape byte-identical};
  \node[inkl, align=center] at (12.6,2.5){\footnotesize 0.1 s $\cdot$ 20 MB\\ 1000$\times$ faster};
  \draw[penf,-{Stealth[length=2mm]}] (4.9,5.0) -- (6.5,5.0);
  \draw[penf,-{Stealth[length=2mm]}] (9.5,5.0) -- (11.1,5.0);
  \platecap{the swiftness was never the language --- only one small thing, done well}
}

% =====================================================================
% 9b. WITHDRAW THE SHARING  (dedup off; colour 80)
% =====================================================================
\art{
  \setsat{80}
  \platetitle{Withdraw the Sharing}
  \def\cy{4.7}
  % ---- left: the fractured casting (C) ----
  \def\lx{4.5}
  \coordinate (TL) at (\lx-0.85,\cy+1.15);\coordinate (TR) at (\lx+0.85,\cy+1.15);
  \coordinate (SR) at (\lx+1.25,\cy+0.5); \coordinate (CU) at (\lx,\cy-1.7);
  \coordinate (SL) at (\lx-1.25,\cy+0.5);
  \begin{scope}\clip (TL)--(TR)--(SR)--(CU)--(SL)--cycle;
     \wfill{wred}{0.5}
     \begin{scope}\clip (\lx-1.3,\cy-1.7) rectangle (\lx+1.3,\cy+0.5);
        \hatch{32}{0.1}{0.4pt}{ink}\hatch{-34}{0.13}{0.4pt}{ink}\end{scope}
  \end{scope}
  \draw[pen2] (TL)--(TR)--(SR)--(CU)--(SL)--cycle;
  \draw[pen] (SL)--(SR);
  \draw[pen2] (\lx-1.0,\cy+0.7) -- (\lx-0.1,\cy+0.2) -- (\lx+0.4,\cy-0.4) -- (\lx+1.1,\cy-0.7);
  \draw[pen] (\lx+0.1,\cy+1.15) -- (\lx-0.15,\cy+0.2) -- (\lx+0.25,\cy-0.8);
  % a shard chipped free, upper right
  \draw[pen2] (\lx+0.85,\cy+1.15) -- (\lx+1.45,\cy+1.5) -- (\lx+1.6,\cy+1.0) -- (\lx+1.25,\cy+0.5);
  \node[ink, lab] at (\lx,\cy-2.35) {\scshape C: it shatters};
  % ---- right: the sound casting (Rust) ----
  \def\rx{11.5}
  \coordinate (tl) at (\rx-0.85,\cy+1.15);\coordinate (tr) at (\rx+0.85,\cy+1.15);
  \coordinate (sr) at (\rx+1.25,\cy+0.5); \coordinate (cu) at (\rx,\cy-1.7);
  \coordinate (sl) at (\rx-1.25,\cy+0.5);
  \begin{scope}\clip (tl)--(tr)--(sr)--(cu)--(sl)--cycle;
     \wfill{wteal}{0.55}
     \begin{scope}\clip (\rx,\cy-1.7) rectangle (\rx+1.3,\cy+0.5);
        \hatch{-28}{0.12}{0.4pt}{inkl}\end{scope}
  \end{scope}
  \draw[pen2] (tl)--(tr)--(sr)--(cu)--(sl)--cycle;
  \draw[pen] (sl)--(sr);
  \draw[pen] (tl)--(cu) (tr)--(cu) (\rx,\cy+1.15)--(\rx,\cy-1.7);
  \fill[paper] (\rx-0.45,\cy+0.72) circle (0.15);
  \node[ink, lab] at (\rx,\cy-2.35) {\scshape Rust: it holds};
  \platecap{switch off the line-reuse --- load-bearing in C, merely cosmetic in Rust}
}

% =====================================================================
% 10. THE BACK-PORT  (colour 84)
% =====================================================================
\art{
  \setsat{84}
  \platetitle{Teaching the Old Engine}
  \foreach \cx/\cy/\r/\lbl in {5.6/4.6/1.5/{C}, 10.4/4.6/1.5/{C\,+\,hash}}{
     \begin{scope}\clip (\cx,\cy) circle (\r);\wfill{wgold}{0.55}\end{scope}
     \foreach \i in {0,...,15}{\pgfmathsetmacro{\a}{\i*22.5}
        \fill[ink] (\cx,\cy)+({\a-4}:\r) -- +({\a-4}:{\r+0.28})
           -- +({\a+4}:{\r+0.28}) -- +({\a+4}:\r) -- cycle;}
     \begin{scope}\clip (\cx,\cy) circle (\r);
        \hatch{30}{0.16}{0.4pt}{ink}\hatch{-34}{0.2}{0.4pt}{inkl}\end{scope}
     \draw[pen2] (\cx,\cy) circle (\r);
     \fill[paper] (\cx,\cy) circle (0.45);\draw[pen] (\cx,\cy) circle (0.45);
     \node[ink] at (\cx,\cy) {\small \lbl};}
  \foreach \a in {0,45,...,315}{\draw[penf] (8.0,4.6)+(\a:0.12) -- +(\a:0.34);}
  \node[ink] at (8.0,2.4) {\scshape\fontsize{24}{26}\selectfont 95$\times$, byte for byte};
  \platecap{put the one trick back into the original itself --- proof it is the idea, not the tongue}
}

% =====================================================================
% 11. THE DOOR UNSEALED  (colour 90)
% =====================================================================
\art{
  \setsat{90}
  \platetitle{The Door Unsealed}
  \begin{scope}\clip (8,3.0) rectangle (10.2,5.4);\wfill{wgold}{0.55}\end{scope}
  \draw[pen2] (8,3.0) rectangle (10.2,5.4);
  \begin{scope}\clip (8,3.0) rectangle (10.2,5.4);
     \clip (9.3,2.8) rectangle (10.2,5.4);
     \hatch{90}{0.12}{0.4pt}{ink}\end{scope}
  \fill[paper] (9.1,4.2) circle (0.2);\draw[pen] (9.1,4.2) circle (0.2);
  \draw[pen2] (9.1,4.2) -- (9.1,3.6);\fill[ink] (9.1,3.6) circle (0.06);
  \draw[pen2, line width=1.4pt] (8.3,5.4) .. controls (8.0,7.0) and (10.0,7.2) .. (10.0,5.9);
  \foreach \x in {3.2,4.6}{
     \draw[pen, inkl] (\x-0.5,3.6) rectangle (\x+0.5,4.7);
     \draw[pen, inkl] (\x-0.3,4.7) arc (180:0:0.3);
     \fill[inkl] (\x,4.15) circle (0.05);}
  \node[inkl] at (3.9,3.1) {\footnotesize frozen};
  \draw[pen,-{Stealth[length=2.4mm]}] (5.6,4.2) -- (7.4,4.2);
  \platecap{the words after the program's name --- once fixed, now free for the solver to choose}
}

% =====================================================================
% 11b. THE SECRET, FOUND BY REASON  (X/Y; colour 92)
% =====================================================================
\art{
  \setsat{92}
  \platetitle{The Secret, Found by Reason}
  \foreach \dx/\ch/\an in {5.2/X/1, 10.8/Y/2}{
     \begin{scope}\clip (\dx,4.8) circle (1.35);
        \wfill{wgold}{0.5}
        \begin{scope}\clip (\dx-1.4,3.45) rectangle (\dx+1.4,4.8);
           \hatch{-28}{0.12}{0.4pt}{inkl}\end{scope}\end{scope}
     \draw[pen2] (\dx,4.8) circle (1.35);
     \draw[penf] (\dx,4.8) circle (1.1);
     \foreach \t in {0,15,...,345}{\draw[pen] ($(\dx,4.8)+(\t:1.1)$) -- ($(\dx,4.8)+(\t:1.32)$);}
     % pointer mark at the top
     \fill[ink] ($(\dx,4.8)+(90:1.46)$) -- ($(\dx,4.8)+(84:1.18)$) -- ($(\dx,4.8)+(96:1.18)$) -- cycle;
     % central disc with the value it landed on
     \fill[paper] (\dx,4.8) circle (0.56);\draw[pen2] (\dx,4.8) circle (0.56);
     \node[ink] at (\dx,4.8) {\fontsize{30}{32}\selectfont \ch};
     \node[ink, lab] at (\dx,2.95) {\scshape arg \an\ reads \ch};}
  % the lock clicks open between the dials
  \foreach \a in {0,30,...,330}{\draw[penf] (8,4.8)+(\a:0.18) -- +(\a:0.5);}
  \node[ink, lab] at (8,3.7) {\footnotesize\itshape it opens};
  \platecap{it crashes only when both land at once --- no keyboard could find it; the solver did}
}

% =====================================================================
% 12. THE IMPARTIAL BALANCE  (colour 95)
% =====================================================================
\art{
  \setsat{95}
  \platetitle{The Impartial Balance}
  \def\cx{8}
  \draw[pen2] (\cx,2.0) -- (\cx,6.3);
  \draw[pen2] (\cx-1.0,2.0) -- (\cx+1.0,2.0);
  \begin{scope}\clip (\cx-1.0,1.7) rectangle (\cx+1.0,2.0);\wfill{wslate}{0.5}\hatch{30}{0.08}{0.4pt}{ink}\end{scope}
  \draw[pen2] (\cx,6.3) -- (\cx-0.28,6.0) -- (\cx+0.28,6.0) -- cycle;
  \draw[pen2] (\cx-3.4,6.25) -- (\cx+3.4,6.25);
  \draw[pen2] (\cx,6.25) circle (0.08);
  \foreach \px/\pig in {-3.4/wwater, 3.4/wamber}{
     \draw[penf] (\cx+\px,6.25) -- (\cx+\px-0.7,4.5);
     \draw[penf] (\cx+\px,6.25) -- (\cx+\px+0.7,4.5);
     \draw[penf] (\cx+\px,6.25) -- (\cx+\px,4.5);
     \begin{scope}\clip (\cx+\px-0.95,4.5) .. controls (\cx+\px-0.6,3.8) and (\cx+\px+0.6,3.8) .. (\cx+\px+0.95,4.5) -- cycle;
        \wfill{\pig}{0.6}\hatch{-25}{0.1}{0.4pt}{inkl}\end{scope}
     \draw[pen2] (\cx+\px-0.95,4.5) .. controls (\cx+\px-0.6,3.8) and (\cx+\px+0.6,3.8) .. (\cx+\px+0.95,4.5);}
  \node[ink] at (\cx-3.4,3.5) {\scshape C rotor};
  \node[ink] at (\cx+3.4,3.5) {\scshape rust rotor};
  \node[ink, lab] at (\cx,4.9) {\scshape\fontsize{30}{32}\selectfont 36 / 36};
  \platecap{same bug, at the same step --- for every benchmark, in both settings, with never a tilt}
}

% =====================================================================
% 13. FROM CIPHER TO PICTURE  (colour 98)
% =====================================================================
\art{
  \setsat{98}
  \platetitle{From Cipher to Picture}
  \begin{scope}\clip (1.8,2.2) rectangle (6.6,6.4);\wfill{wsky}{0.35}\end{scope}
  \draw[pen] (1.8,2.2) rectangle (6.6,6.4);
  \foreach \ly in {2.6,3.0,...,6.0}{
     \pgfmathsetmacro{\lw}{2.4+rnd*1.8}
     \draw[penf, decorate, decoration={random steps, segment length=3pt, amplitude=0.5pt}]
        (2.2,\ly) -- (2.2+\lw,\ly);}
  \node[ink] at (4.2,6.9) {\footnotesize\itshape the solver's raw witness};
  \draw[pen2,-{Stealth[length=2.8mm]}] (7.0,4.3) -- (8.8,4.3);
  \coordinate (g0) at (12.4,6.2);
  \coordinate (g1) at (10.6,4.6);\coordinate (g2) at (14.2,4.6);
  \coordinate (g3) at (9.8,3.0);\coordinate (g4) at (11.6,3.0);
  \draw[pen] (g0)--(g1) (g0)--(g2) (g1)--(g3) (g1)--(g4);
  \begin{scope}\clip (g0) circle (0.42);\wfill{wred}{0.7}\end{scope}
  \foreach \a in {0,45,...,315}{\draw[penf] (g0)+(\a:0.5) -- +(\a:0.74);}
  \draw[pen2] (g0) circle (0.42);
  \node[ink] at (g0) {\scriptsize\scshape bad};
  \foreach \g/\pig in {g1/wteal, g2/wamber, g3/wviolet, g4/wgreen}{
     \begin{scope}\clip (\g) circle (0.34);\wfill{\pig}{0.65}\end{scope}
     \begin{scope}\clip (\g) circle (0.34);\hatch{40}{0.1}{0.35pt}{inkl}\end{scope}
     \draw[pen] (\g) circle (0.34);}
  \platecap{a wall of binary becomes a graph you can point at}
}

% =====================================================================
% 14. THE SUMMIT  -- fully tinted (colour 100)
% =====================================================================
\art{
  \setsat{100}
  \draw[pen] (0.7,0.7) rectangle (15.3,8.3);
  \draw[penf] (0.95,0.95) rectangle (15.05,8.05);
  % sky wash
  \begin{scope}\clip (0.95,2.6) rectangle (15.05,8.05);\wfill{wsky}{0.5}\end{scope}
  % foreground meadow
  \begin{scope}\clip (0.95,0.95) rectangle (15.05,2.6);\wfill{wgreen2}{0.55}\end{scope}
  % mountain: green wash, then ink shadow hatching
  \begin{scope}\clip (3.0,2.6) -- (8.0,7.4) -- (12.6,2.6) -- cycle;\wfill{wgreen}{0.6}\end{scope}
  \begin{scope}
    \clip (4.6,2.6) -- (8.0,7.4) -- (8.0,2.6) -- cycle;
    \wfill{wslate}{0.45}
    \hatch{60}{0.1}{0.4pt}{ink}
  \end{scope}
  \draw[pen2] (3.0,2.6) -- (8.0,7.4) -- (12.6,2.6);
  \draw[pen] (8.0,7.4) -- (6.6,5.0) -- (7.3,4.2) -- (6.2,3.3);
  % clouds (white, masking the sky)
  \foreach \cx/\cy in {3.4/7.2, 11.8/7.3}{
     \fill[paper, opacity=0.9] (\cx,\cy) .. controls (\cx-0.8,\cy+0.5) and (\cx-1.4,\cy-0.1) .. (\cx-1.7,\cy)
        .. controls (\cx-1.2,\cy-0.3) and (\cx+1.2,\cy-0.3) .. (\cx+1.7,\cy)
        .. controls (\cx+1.2,\cy+0.4) and (\cx+0.6,\cy+0.4) .. (\cx,\cy);
     \draw[penf, inkl] (\cx,\cy) .. controls (\cx-0.8,\cy+0.5) and (\cx-1.4,\cy-0.1) .. (\cx-1.7,\cy)
        .. controls (\cx-1.2,\cy-0.3) and (\cx+1.2,\cy-0.3) .. (\cx+1.7,\cy)
        .. controls (\cx+1.2,\cy+0.4) and (\cx+0.6,\cy+0.4) .. (\cx,\cy);}
  % flag (red)
  \draw[pen2] (8.0,7.4) -- (8.0,8.0);
  \fill[wred] (8.0,8.0) -- (8.55,7.85) -- (8.0,7.7) -- cycle;
  \draw[pen2] (8.0,8.0) -- (8.55,7.85) -- (8.0,7.7);
  % a parchment cartouche so the closing words sit clear of the painted slopes
  \fill[paper, opacity=0.9, rounded corners=5pt] (2.7,0.8) rectangle (13.3,2.55);
  \draw[penf, rounded corners=5pt] (2.7,0.8) rectangle (13.3,2.55);
  \node[ink] at (8,2.05) {\scshape\fontsize{18}{20}\selectfont proving it still tells the truth};
  \node[ink] at (8,1.5) {\itshape\footnotesize was the whole climb. \quad thank you.};
  \node[inkl] at (8,1.05) {\scriptsize github.com/jaspek/rotor-rust \,\textbar\, rotor-c-hashcons \,\textbar\, jaspek.github.io/rotor-rust};
}

\end{document}
